\section{確率的プログラミングの応用1;項目反応理論}
\subsection{授業内容}
\subsubsection{科目の中でのこのコマの位置づけ}
ここまでで線一般形モデル，一般化線形モデル，階層線形モデル，混合分布モデルと定型的な分析モデルについて一通り学んできた。
以後はアラカルト的に，確率的プログラミングの応用による柔軟なモデリング例のトピックスを取り上げる。
最初に扱うのは，既に前期に学んだ項目反応理論のモデリングである。尺度作成の文脈で，理論的概要は一通り説明が終わっているところであるが，改めて確認するとともに確率的モデリングとして実装する。
確率モデルとして考えると，0/1の反応に対するロジスティック回帰の応用であり，実装自体は既有知識の応用で可能であろう。コーディングのポイントとして，long型データ(tidy data)にしておくことで欠損値が含まれる場合も対応できるようになることが挙げられる。
\subsubsection{コマ主題細目}
\begin{description}
	\item[ロジスティックモデルの復習] 本講では前期のうちに，ロジスティックモデルについての理論的説明と，Rの\verb|ltm|パッケージによる実践例を解説済みである。とはいえ，以前学んでから時間が空いているので，あるいは後期のみ履修する学生もいることが考えられるんので，授業の冒頭で15分程度の時間をかけて，大まかな理論・モデルの復習をしておく必要があるだろう。
	\item[ロジスティック回帰モデルでの実装] ロジスティック回帰分析を思い出しつつ，1PL,2PL,3PLモデルそれぞれを\verb|transformed parameters|ブロックで記述することを演習で学ぶ。
	\item[整然データでの分析] データを整然データの形にして分析することで，欠損値が含まれないデータセットを作って分析に応用することができる。ここでは個人と項目それぞれを識別する変数が必要になるが，これまで学んできた技術で十分対応可能であると考えられる。
\end{description}
\subsubsection{キーワード}
\begin{itemize}
	\item 項目反応理論
	\item 1PLロジスティックモデル
	\item 2PLロジスティックモデル
	\item 3PLロジスティックモデル
	\item 整然データ
\end{itemize}
\subsection{授業情報}
\paragraph{コマの展開方法}
講義/遠隔可/演習
\subsubsection{予習・復習課題}
\paragraph{予習} これまでの知識や技術を組み合わせて問題に対応することになる。項目反応理論とロジスティックモデル，GLMにおけるロジスティック回帰分析，データハンドリングにおける整然データの考え方など，これまでの資料に戻って復習しておくと良い。
\paragraph{復習} 自分で描いたモデルがRのパッケージが出す答えとどの程度一致するのかを確認しておこう。また欠損値がある被験者の被験者母数は，その確信区間が広くなると考えられる。なぜそうなるかを改めて考え，実際のデータ適用例で確認しておこう。